\documentclass[11pt,letterpaper]{article}

% ==========================================================
%  RiskGate - Documentación técnica completa (español)
% ==========================================================

\usepackage[spanish,es-noshorthands]{babel}
\usepackage[utf8]{inputenc}
\usepackage[T1]{fontenc}
\usepackage{lmodern}
\usepackage[margin=1in]{geometry}
\usepackage{microtype}
\usepackage{amsmath}
\usepackage{amssymb}
\usepackage{array}
\usepackage{booktabs}
\usepackage{longtable}
\usepackage{enumitem}
\usepackage{xcolor}
\usepackage{listings}
\usepackage{hyperref}

\hypersetup{
  colorlinks=true,
  linkcolor=blue,
  urlcolor=blue,
  citecolor=blue,
  pdftitle={RiskGate - Documentacion tecnica completa},
  pdfauthor={RiskGate}
}

\definecolor{rgcode}{gray}{0.96}
\lstdefinestyle{riskgate}{
  basicstyle=\ttfamily\small,
  breaklines=true,
  frame=single,
  columns=fullflexible,
  backgroundcolor=\color{rgcode},
  keepspaces=true,
  showstringspaces=false
}
\lstset{style=riskgate}

\newcommand{\RiskGate}{\textsc{RiskGate}}
\newcommand{\MXN}{\ensuremath{\mathrm{MXN}}}
\newcommand{\USD}{\ensuremath{\mathrm{USD}}}
\newcommand{\ind}[1]{\mathbb{1}\{#1\}}
\newcommand{\clp}{\operatorname{clamp}}
\setlist[itemize]{leftmargin=1.3em}
\setlist[enumerate]{leftmargin=1.5em}

\title{\textbf{\RiskGate}\\[2mm]
\large Documentación técnica completa del proyecto\\
y recomendaciones de evolución}
\author{Proyecto \RiskGate{} --- MVP local-first}
\date{Mayo de 2026}

\begin{document}

\maketitle

\begin{abstract}
\noindent
\RiskGate{} es una aplicación \emph{local-first} de gestión de portafolio y de
riesgo en opciones. Su función no es ejecutar órdenes, predecir precios ni dar
asesoría financiera, sino \textbf{imponer disciplina de proceso} antes de
entrar a una operación: obliga a documentar cada idea de trade, la evalúa con un
modelo cuantitativo transparente, rechaza estructuras incompletas o ilíquidas y
asigna un tamaño de riesgo conservador acotado por los límites de la cuenta.
Este documento describe de forma exhaustiva la arquitectura, el modelo de datos,
las pantallas, el motor de evaluación (puntaje ponderado adaptativo, calidad de
opciones continua, indicadores técnicos, analítica probabilística, simulación
Monte Carlo, sentimiento de Reddit y concentración sectorial), el bot de
investigación en Python y un conjunto detallado de recomendaciones y
modificaciones de \emph{proxies} para hacer el sistema más robusto.
\end{abstract}

\tableofcontents
\newpage

% ==========================================================
\section{Propósito y filosofía del sistema}
% ==========================================================

\RiskGate{} fue diseñado como un sistema de apoyo a la decisión para trading
manual de acciones y opciones. La regla central que gobierna todo el diseño es:

\begin{quote}
\itshape
Ninguna operación de opciones debe aprobarse a menos que el plan escrito, el
puntaje ponderado, los controles de liquidez, las reglas de vencimiento y los
topes de riesgo del portafolio se cumplan simultáneamente.
\end{quote}

Los objetivos concretos del sistema son:

\begin{itemize}[noitemsep]
  \item mantener varios portafolios (cuentas) aislados entre sí;
  \item importar datos de portafolio exportados manualmente desde GBM;
  \item registrar posiciones, efectivo, exposición y asignación de activos;
  \item mantener una \emph{watchlist} de tickers seguidos por portafolio;
  \item obligar a que cada idea de opciones tenga tesis escrita, nivel de
        invalidación, pérdida máxima y plan de salida;
  \item puntuar ideas combinando contexto de mercado, técnico, calidad del
        contrato y calidad del plan;
  \item rechazar configuraciones incompletas o inseguras;
  \item calcular un tamaño de riesgo sugerido acotado por los límites de cuenta; y
  \item registrar operaciones entradas y cerradas en una bitácora para revisión.
\end{itemize}

\subsection{Principios de seguridad (local-first)}

El diseño local-first se eligió porque la aplicación maneja información
financiera. El MVP evita deliberadamente:

\begin{itemize}[noitemsep]
  \item almacenar contraseñas de GBM o del broker;
  \item hacer \emph{scraping} de páginas de inicio de sesión;
  \item sincronizar datos privados a la nube;
  \item trading automático o colocación de órdenes.
\end{itemize}

La base de datos SQLite vive en la máquina del usuario y todas las fuentes
externas se consumen en modo \emph{solo lectura}.

% ==========================================================
\section{Arquitectura y stack tecnológico}
% ==========================================================

El producto se compone de dos piezas independientes:

\begin{enumerate}[noitemsep]
  \item \textbf{La aplicación web} (Next.js / React / TypeScript), que es el MVP
        principal y contiene el motor de evaluación.
  \item \textbf{El bot de investigación} (\texttt{options-strategy-bot}, Python),
        un banco de pruebas independiente para escanear, puntuar, hacer
        \emph{backtest} y llevar un libro de \emph{paper trading}.
\end{enumerate}

\begin{longtable}{p{0.30\linewidth} p{0.62\linewidth}}
\toprule
\textbf{Capa} & \textbf{Implementación} \\
\midrule
Frontend & Next.js (App Router), React, TypeScript \\
Estilos & Tailwind CSS, \texttt{@tailwindcss/forms} \\
Lógica de servidor & Server Actions de Next.js (\texttt{app/actions.ts}) \\
Base de datos & SQLite \\
ORM & Prisma Client (\texttt{@prisma/client} v6) \\
Gráficas & Recharts \\
Iconos & Lucide React \\
Validación & Zod \\
Importación Excel & paquete \texttt{read-excel-file} \\
Datos de mercado & Yahoo Finance (endpoint público de \emph{chart}) \\
Cadena de opciones & Alpaca Options Data API (solo lectura, opcional) \\
Bot de investigación & Python 3, Black--Scholes propio, Yahoo/Alpaca/CSV \\
Runtime & Servidor de desarrollo local de Node.js \\
\bottomrule
\end{longtable}

\noindent
Todas las páginas se sirven con \texttt{dynamic = "force-dynamic"} para reflejar
datos en tiempo real. Cada operación de mutación se ejecuta como \emph{server
action} y revalida las rutas afectadas tras escribir en la base de datos.

% ==========================================================
\section{Modelo de datos}
% ==========================================================

El esquema Prisma (SQLite) organiza la información alrededor de la entidad
\texttt{Portfolio}, que aísla cuentas. Las entidades principales son:

\begin{longtable}{p{0.24\linewidth} p{0.68\linewidth}}
\toprule
\textbf{Modelo} & \textbf{Propósito} \\
\midrule
\texttt{Portfolio} & Cuenta o portafolio independiente; ancla del aislamiento de datos. \\
\texttt{Position} & Tenencias manuales o importadas (acción, ETF, opción, efectivo). \\
\texttt{WatchlistItem} & Tickers seguidos con cotización, cambio diario y fuente. \\
\texttt{TradeIdea} & Idea de opciones pre-entrada: inputs, puntajes, decisión, riesgo sugerido, razones y advertencias. \\
\texttt{OptionLeg} & Pierna individual de una estructura multi-leg (tipo, lado, strike, vencimiento, Greeks, IV, OI, volumen). \\
\texttt{ModelScore} & Registro normalizado de puntaje, uno a uno con la idea. \\
\texttt{TradeScoreSnapshot} & Foto diaria del puntaje de una idea para seguir su evolución. \\
\texttt{RiskSettings} & Límites de riesgo y parámetros de dimensionamiento por portafolio. \\
\texttt{TradeJournal} & Operaciones entradas y cerradas, ligadas a una idea. \\
\texttt{PriceBar} & Barras OHLCV históricas por símbolo y timeframe. \\
\texttt{PortfolioSnapshot} & Foto puntual del valor del portafolio (intradía e histórico). \\
\texttt{RedditDocument} & Documentos de Reddit normalizados para análisis de sentimiento. \\
\bottomrule
\end{longtable}

\subsection{Aislamiento por portafolio}

Los modelos \texttt{Position}, \texttt{TradeIdea}, \texttt{RiskSettings},
\texttt{WatchlistItem}, \texttt{PriceBar} y \texttt{PortfolioSnapshot} llevan un
campo \texttt{portfolioId}. El portafolio activo se transmite por la URL:

\begin{lstlisting}
?portfolio=<portfolioId>
\end{lstlisting}

Las entradas de bitácora se ligan a través de la idea de trade, por lo que están
indirectamente aisladas por portafolio.

% ==========================================================
\section{Módulos y pantallas de la aplicación}
% ==========================================================

La barra lateral (\texttt{components/sidebar.tsx}) contiene el selector de
portafolio y la navegación a diez pantallas.

\subsection{Dashboard}
Resume el portafolio seleccionado: equity total, efectivo, exposición a
acciones/ETF y opciones, cambio intradía, riesgo de opciones utilizado, tope de
pérdida mensual, gráfica de asignación (pastel), curva de equity, desempeño
intradía, gráfica de uso de riesgo y la tabla de ideas recientes con su decisión.

\subsection{Portafolio}
Permite crear, editar y eliminar posiciones e importar el Excel de GBM. Cada
posición tiene símbolo, tipo de activo, mercado (MX/US), cantidad, costo
promedio, precio actual, moneda, cuenta y notas. Refresca precios desde Yahoo y
registra \emph{snapshots} para el desempeño intradía.

\subsection{Importación de Excel de GBM}
El importador lee el libro de detalle de GBM, ignora encabezados de sección
(p.\,ej.\ \texttt{Mercado de Capitales USA}, \texttt{Liquidez}), interpreta filas
de tickers, convierte \texttt{efectivo} en una posición \texttt{CASH}, reconoce
ETFs comunes y, opcionalmente, convierte valores USD a MXN con el tipo de cambio
ingresado. Puede reemplazar todas las posiciones o fusionar por símbolo.

\subsection{Watchlist}
Almacena tickers seguidos por portafolio (símbolo, nombre, mercado, moneda,
fuente, precio, cierre previo, cambio diario, marca de tiempo). Permite refrescar
cotizaciones desde Yahoo o editarlas manualmente. Es la base para el cálculo
automático de indicadores técnicos.

\subsection{Discover}
Pantalla de descubrimiento: escanea el mercado y la watchlist para encontrar
candidatos. Filtros por perfil (\emph{quality} o \emph{movers}), puntaje mínimo,
volumen mínimo e inclusión de Reddit. Muestra puntaje de descubrimiento, sesgo
(alcista/bajista), volatilidad realizada, menciones en Reddit y probabilidad de
liquidez de opciones, con botones para añadir a watchlist o crear una idea.

\subsection{Options Chain}
Navega la cadena de opciones de Alpaca por subyacente, vencimiento y strike.
Muestra bid/ask/mid, delta, gamma, theta y vega, y permite pre-llenar el
formulario de idea con el contrato elegido.

\subsection{Trade Ideas}
Es el \emph{núcleo de la compuerta de riesgo}. Cada posible operación se escribe
aquí antes de entrar. El formulario captura símbolo, dirección, estrategia,
tesis, catalizador, nivel de invalidación, objetivo, vencimiento, prima, pérdida
máxima, recompensa esperada, plan de salida, estructura multi-leg, contexto de
mercado, señales técnicas e indicador de \emph{chasing}. Por cada idea muestra
los cuatro puntajes de módulo, la decisión, el riesgo sugerido, la analítica de
opciones (PoP, EV, Kelly, VaR), el histograma Monte Carlo, el diagrama de payoff,
el desglose técnico y de calidad de opciones, el sentimiento de Reddit y las
listas de razones y advertencias.

\subsection{Risk Model}
Tablero del perfil de riesgo: riesgo base por trade, límite del \emph{sleeve} de
opciones, pesos por régimen, bandas de decisión, multiplicadores, Greeks
agregados del portafolio y analítica de concentración.

\subsection{Model Guide}
Explicación didáctica del modelo dentro de la app: módulos de puntaje, hard
stops, topes de aprobación, bandas de decisión, fórmula de dimensionamiento,
metodología de sentimiento y hoja de ruta de robustez.

\subsection{Journal}
Bitácora de operaciones: enlaza la idea, fechas y precios de entrada/salida, P\&L
realizado, resultado, etiquetas de error y notas. Calcula \emph{win rate},
ganancia/pérdida promedio, \emph{expectancy}, proxies de Sharpe y Sortino,
duración media e intervalos de confianza por \emph{bootstrap}.

\subsection{Settings}
Gestión de portafolios (crear/editar/eliminar) y de los límites de riesgo:
valor total, \emph{sleeve} de opciones, riesgo base por trade, tope de pérdida
mensual, tope de riesgo abierto y la configuración de sentimiento de Reddit.

% ==========================================================
\section{Pipeline de datos de mercado}
% ==========================================================

\subsection{Yahoo Finance (\texttt{lib/market-data.ts})}
Obtiene barras OHLCV (por defecto 6 meses diarios) y cotizaciones. A partir del
historial construye un \emph{resumen técnico} por símbolo y el
\texttt{MarketContext} agregado:

\begin{itemize}[noitemsep]
  \item \textbf{Tendencia de régimen}: \emph{uptrend}/\emph{downtrend}/\emph{neutral}
        según precio vs.\ SMA20/SMA50 y cambio a 20 días.
  \item \textbf{Condición de VIX}: \emph{low} ($<15$), \emph{normal},
        \emph{elevated} ($22$--$30$), \emph{spiking} ($\ge 30$ o salto diario
        $\ge +15\%$).
  \item \textbf{Amplitud (breadth)}: proporción de símbolos con precio
        $\ge$ SMA50. \emph{Strong} si $\ge 65\%$, \emph{weak} si $\le 40\%$.
  \item \textbf{Riesgo macro}: \emph{high} si el VIX está \emph{spiking};
        \emph{low} si VIX bajo y breadth fuerte; \emph{medium} en otro caso.
\end{itemize}

\subsection{Alpaca Options Data API (\texttt{lib/alpaca-options.ts})}
Consulta \emph{snapshots} de la cadena de opciones (Greeks, IV, bid/ask, volumen,
open interest) en modo solo lectura. Decodifica el símbolo OCC
\texttt{[ROOT][YY][MM][DD][C/P][STRIKE$\times$1000]} y normaliza respuestas en
múltiples formatos. Requiere \texttt{APCA\_API\_KEY\_ID} y
\texttt{APCA\_API\_SECRET\_KEY}; si no están configuradas, la app opera con datos
manuales.

% ==========================================================
\section{Motor de evaluación de ideas de trade}
% ==========================================================

El corazón del sistema es \texttt{evaluateTradeIdea()} en \texttt{lib/risk.ts}.
A diferencia de versiones anteriores con ajustes discretos, el modelo actual usa
\textbf{funciones de calidad continuas}, \textbf{pesos adaptativos al régimen} y
\textbf{capas de superposición} probabilísticas.

\subsection{Estructura general}

El puntaje total se compone de cuatro módulos en escala $[0,100]$:

\begin{longtable}{p{0.24\linewidth} p{0.16\linewidth} p{0.50\linewidth}}
\toprule
\textbf{Módulo} & \textbf{Peso típico} & \textbf{Propósito} \\
\midrule
Régimen de mercado ($M$) & 25--35\% & Contexto amplio: SPY, QQQ, VIX, breadth, macro. \\
Técnico ($T$) & 20--30\% & Alineación direccional con tendencia, medias, RSI, volumen, S/R. \\
Calidad de opciones ($O$) & 25--35\% & Liquidez y estructura del contrato. \\
Calidad del trade ($Q$) & 15--20\% & Disciplina del plan escrito. \\
\bottomrule
\end{longtable}

La función de recorte es:
\begin{equation}
\clp(x,0,100) = \min(100,\max(0,x)).
\end{equation}

\subsection{Pesos adaptativos al régimen}

Los pesos $w$ se ajustan según el contexto (\texttt{getAdaptiveWeights}):

\begin{longtable}{p{0.34\linewidth} cccc}
\toprule
\textbf{Régimen} & $w_M$ & $w_T$ & $w_O$ & $w_Q$ \\
\midrule
VIX elevado/spiking o SPY bajista & 0.30 & 0.20 & 0.35 & 0.15 \\
SPY alcista & 0.25 & 0.30 & 0.25 & 0.20 \\
Neutral (por defecto) & 0.35 & 0.25 & 0.25 & 0.15 \\
\bottomrule
\end{longtable}

El puntaje base ponderado es:
\begin{equation}
S_{\text{base}} = w_M M + w_T T + w_O O + w_Q Q.
\end{equation}

\subsection{Puntaje de régimen de mercado}

Parte de 50 y aplica ajustes discretos:
\begin{align}
M = \clp\big(&50
 + 15\ind{\text{SPY alcista}} - 15\ind{\text{SPY bajista}} \notag\\
&+ 15\ind{\text{QQQ alcista}} - 15\ind{\text{QQQ bajista}} \notag\\
&+ 10\ind{\text{breadth fuerte}} - 10\ind{\text{breadth débil}} \notag\\
&+ 5\ind{\text{VIX bajo}} - 10\ind{\text{VIX elevado}} - 25\ind{\text{VIX spiking}} \notag\\
&+ 5\ind{\text{macro bajo}} - 20\ind{\text{macro alto}},\,0,100\big).
\end{align}

\subsection{Puntaje técnico}

Cuando hay barras OHLCV disponibles, el técnico se calcula con indicadores reales
(\texttt{lib/indicators.ts}) en lugar de etiquetas manuales:

\begin{itemize}[noitemsep]
  \item \textbf{SMA}$_n$: media simple de los últimos $n$ cierres.
  \item \textbf{EMA}$_n$: media exponencial con multiplicador $\alpha = 2/(n+1)$.
  \item \textbf{RSI}$_{14}$: $\,RSI = 100 - \dfrac{100}{1+RS},\ RS = \dfrac{\overline{\text{ganancia}}}{\overline{\text{pérdida}}}$.
  \item \textbf{ATR}$_{14}$: media del \emph{true range}
        $\max(H-L,\,|H-C_{-1}|,\,|L-C_{-1}|)$.
  \item \textbf{Volumen relativo}: volumen actual / promedio base.
  \item \textbf{BWAP}: $\sum (p_{\text{típico}}\cdot V)/\sum V$, con
        $p_{\text{típico}}=(H+L+C)/3$.
\end{itemize}

El puntaje técnico continuo pondera siete sub-calidades:
trend ($25\%$), alineación de medias ($20\%$), RSI ($15\%$), volumen ($15\%$),
BWAP ($10\%$), soporte/resistencia ($10\%$) y riesgo por ATR ($5\%$). Si no hay
snapshot calculado, el sistema recurre a un modelo de respaldo con ajustes
discretos sobre una base de 50.

\subsection{Puntaje de calidad de opciones (continuo)}

En \texttt{lib/risk/optionsQuality.ts}, cada dimensión usa una función suave en
$[0,1]$ que luego se pondera a $[0,100]$:

\begin{align}
\text{spread}(s) &= \frac{1}{1+\left(\frac{s-2}{5}\right)^3}, \quad s>2\%, \\[1mm]
\text{volumen}(v) &= \clp\!\left(\sqrt{v/500}\right), \\
\text{OI}(oi) &= \clp\!\left(\sqrt{oi/1000}\right), \\
\text{ivRank}(r) &= 1 - (r/100)^2, \\
\text{theta}(\theta, m) &= 1 - \left|\theta/m\right|,
\end{align}
y la calidad de DTE tiene su máximo en $21$--$60$ días y cae a cero por debajo de
5 o por encima de 240 días. Los pesos son: spread $30\%$, DTE $20\%$, volumen
$15\%$, open interest $15\%$, IV rank $8\%$, tamaño de prima $7\%$ y theta $5\%$.

\subsection{Puntaje de calidad del trade (continuo)}

\begin{equation}
Q = 100\big(0.25\,q_{\text{tesis}} + 0.10\,q_{\text{cat}} + 0.15\,q_{\text{inval}}
+ 0.15\,q_{\text{salida}} + 0.20\,q_{R} + 0.15\,q_{\text{chase}}\big),
\end{equation}
donde $q_{\text{tesis}} = \clp(\text{long}/80)$ (longitud del texto),
$q_{\text{cat}} = 1$ si hay catalizador (si no, $0.35$),
$q_{\text{inval}} = 1$ si existe nivel de invalidación,
$q_{\text{salida}} = \clp(\text{long}/80)$,
$q_R = \clp(R/3)$ con $R = \text{recompensa}/\text{pérdida máxima}$, y
$q_{\text{chase}} = 0$ si la idea está marcada como \emph{chasing}.

\subsection{Capas de superposición}

Sobre $S_{\text{base}}$ se aplican ajustes aditivos:

\begin{itemize}[noitemsep]
  \item \textbf{Monte Carlo}: $+5$ si EV$>0$ y PoP$>55\%$; $-5$ si EV$<0$;
        $-10$ si CVaR$_{95}\ge 0.85\cdot$ pérdida máxima.
  \item \textbf{Sentimiento de Reddit} (si está activo): $+3$ si concuerda con la
        dirección y la confianza $>0.65$; $-3$ si se opone; penalizaciones por
        \emph{hype} o por preocupaciones severas.
\end{itemize}

\begin{equation}
S_{\text{total}} = \clp\big(S_{\text{base}} + \Delta_{\text{MC}} + \Delta_{\text{Reddit}},\,0,100\big).
\end{equation}

% ==========================================================
\section{Analítica de opciones}
% ==========================================================

El módulo \texttt{lib/options-math.ts} provee analítica probabilística.

\subsection{Probabilidad de ganancia (PoP)}
Bajo un modelo lognormal del subyacente, para un strike $K$, precio $S_0$,
volatilidad implícita $\sigma$ y tiempo a vencimiento $\tau$ (en años):
\begin{equation}
d = \frac{\ln(K/S_0) + \tfrac{1}{2}\sigma^2\tau}{\sigma\sqrt{\tau}}, \qquad
\text{PoP}_{\text{call}} = 1 - N(d), \quad \text{PoP}_{\text{put}} = N(d),
\end{equation}
donde $N(\cdot)$ es la CDF normal estándar (vía función de error). El punto de
equilibrio es $K + \text{prima}$ (call) o $K - \text{prima}$ (put).

\subsection{Valor esperado y Kelly}
\begin{align}
\text{EV} &= \text{PoP}\cdot\text{recompensa} - (1-\text{PoP})\cdot\text{pérdida máxima}, \\
f^{*} &= \max\!\left(0,\ \frac{\text{PoP}\cdot b - (1-\text{PoP})}{b}\right),
\end{align}
con $b$ el cociente recompensa/riesgo. El sistema usa \textbf{Kelly fraccional al
$25\%$} ($f_{1/4} = 0.25\,f^{*}$) como tope de tamaño.

\subsection{VaR y CVaR (un día)}
Usando una aproximación delta--gamma--vega más decaimiento por theta:
\begin{align}
\Delta S &= S_0 \cdot \frac{\sigma}{\sqrt{252}}, \\
\sigma_{\text{proxy}} &= |\delta|\,\Delta S + \tfrac{1}{2}|\gamma|\,\Delta S^2 + |\nu|\,\sigma\cdot 0.1, \\
\text{VaR}_{95} &= \sigma_{\text{proxy}}\,z_{95} + \theta_{\text{loss}}, \quad z_{95}=1.645, \\
\text{CVaR}_{95} &= \sigma_{\text{proxy}}\cdot \frac{\phi(z_{95})}{0.05} + \theta_{\text{loss}}.
\end{align}

\subsection{Greeks en dólares y agregación}
La exposición se escala por $100$ acciones por contrato. La función
\texttt{aggregateOptionAnalytics} suma delta, gamma, theta y vega de todas las
ideas aprobadas/entradas para mostrar los Greeks del portafolio.

\subsection{Estructuras multi-leg (\texttt{lib/options/multileg.ts})}
El payoff al vencimiento por pierna es $\max(S_T-K,0)$ (call) o $\max(K-S_T,0)$
(put). El P\&L de una pierna long es $q\cdot 100\cdot(\text{payoff}-\text{prima})$
y se invierte para piernas short. La calculadora evalúa el payoff en una malla de
puntos (101 regulares más los strikes y el precio actual), localiza los puntos de
equilibrio por interpolación lineal y deriva débito/crédito neto, pérdida y
ganancia máximas, y Greeks netos. Existen plantillas para \emph{long call/put},
\emph{bull call} y \emph{bear put debit spreads}, \emph{protective put} y
\emph{collar}, y un detector de piernas short desnudas que actúa como hard stop.

% ==========================================================
\section{Simulación Monte Carlo}
% ==========================================================

En \texttt{lib/options/monteCarlo.ts} se simulan trayectorias del subyacente con
un \emph{movimiento browniano geométrico} (GBM):
\begin{equation}
S_T = S_0 \exp\!\left[\left(\mu - q - \tfrac{1}{2}\sigma^2\right)\tau + \sigma\sqrt{\tau}\,Z\right],
\quad Z\sim\mathcal{N}(0,1),
\end{equation}
con deriva $\mu$ (por defecto 0) y rendimiento por dividendos $q$. Para cada
trayectoria se valúa la estructura multi-leg y se construye la distribución de
P\&L, de la que se obtienen PoP, valor esperado, mediana, percentiles P5/P95,
CVaR$_{95}$ y un histograma. El motor usa una semilla determinista (función del
símbolo y el puntaje) para reproducibilidad, con miles de trayectorias por
evaluación.

% ==========================================================
\section{Sentimiento de Reddit (opcional)}
% ==========================================================

El subsistema \texttt{lib/sentiment/reddit/} resume narrativa pública como señal
cualitativa, siempre opcional y desactivada por defecto.

\begin{itemize}[noitemsep]
  \item \textbf{Recolección}: busca en Reddit vía OAuth posts sobre el ticker y
        sus catalizadores, ponderando subreddits (options $1.0$, investing/stocks
        $0.9$, wallstreetbets $0.65$).
  \item \textbf{Limpieza}: elimina URLs y menciones, decodifica HTML, colapsa
        espacios y tokeniza.
  \item \textbf{Embeddings TF--IDF locales}: sin modelos externos; se calcula
        similitud por coseno $\cos(A,B)=\frac{A\cdot B}{\lVert A\rVert\,\lVert B\rVert}$.
  \item \textbf{Recuperación}: puntaje $= 0.65\,\text{coseno} + 0.20\,\text{recencia} + 0.10\,\text{engagement} + 0.05\,\text{peso de subreddit}$.
  \item \textbf{Clasificación}: produce etiquetas (\emph{strong\_bullish},
        \emph{bullish}, \emph{mixed}, \emph{bearish}, \emph{strong\_bearish},
        \emph{hype}, \emph{fear}, \emph{low\_signal}) a partir de puntajes de
        sentimiento, hype y miedo, junto con argumentos, preocupaciones
        recurrentes y catalizadores mencionados.
\end{itemize}

\noindent
El sentimiento ajusta tanto el puntaje (capa de superposición) como un
multiplicador de riesgo ($0.75\times$ ante hype alto o desacuerdo con alta
confianza).

% ==========================================================
\section{Concentración sectorial}
% ==========================================================

En \texttt{lib/concentration.ts} se mide la concentración con el índice de
Herfindahl--Hirschman:
\begin{equation}
\text{HHI} = \sum_k w_k^2, \qquad w_k = \frac{\text{riesgo}_k}{\text{riesgo total}},
\end{equation}
donde $k$ recorre sectores. El multiplicador de penalización es $0.5$ si
$\text{HHI}>0.5$ (concentración alta), $0.75$ si $\text{HHI}>0.35$ (moderada) y
$1.0$ en otro caso. Este multiplicador entra directamente en el dimensionamiento.

% ==========================================================
\section{Decisión y dimensionamiento de riesgo}
% ==========================================================

\subsection{Bandas de decisión}
\begin{longtable}{p{0.30\linewidth} p{0.35\linewidth}}
\toprule
\textbf{Puntaje total} & \textbf{Decisión inicial} \\
\midrule
$S<60$ & reject \\
$60\le S<75$ & watchlist \\
$75\le S<85$ & approved\_small \\
$S\ge 85$ & approved\_normal \\
\bottomrule
\end{longtable}

\subsection{Hard stops (rechazo forzado)}
Sin importar el puntaje, cualquiera de estos fuerza \emph{reject}:
\begin{itemize}[noitemsep]
  \item falta tesis escrita, nivel de invalidación, pérdida máxima o plan de salida;
  \item spread bid/ask $>15\%$; DTE $<7$; prima $>1\%$ del portafolio;
  \item volumen $<25$ y open interest $<100$ simultáneamente;
  \item datos de pierna faltantes que impiden calcular la pérdida máxima;
  \item estrategia de riesgo definido sin pérdida máxima determinable;
  \item piernas short desnudas.
\end{itemize}

\subsection{Topes de aprobación (limitan a watchlist)}
La idea no puede superar \emph{watchlist} si: está marcada como \emph{chasing},
IV rank $>80$, IV del contrato $>85\%$, theta $>8\%$ del mid por día, o gamma alto
con DTE $<30$.

\subsection{Fórmula de dimensionamiento}

El riesgo por multiplicadores es:
\begin{equation}
R_{\text{mult}} = B \cdot C(S)\cdot L(s)\cdot G\cdot V(r)\cdot K_{\text{conc}}\cdot K_{\text{MC}}\cdot K_{\text{sent}},
\end{equation}
con $B$ el riesgo base por trade (\MXN), y los multiplicadores:

\begin{longtable}{p{0.30\linewidth} p{0.62\linewidth}}
\toprule
\textbf{Multiplicador} & \textbf{Regla} \\
\midrule
Confianza $C(S)$ & $0$ si $S<75$; $0.25$ ($75$--$80$); $0.50$ ($80$--$85$); $0.75$ ($85$--$90$); $1.0$ ($\ge 90$). \\
Liquidez $L(s)$ & $0$ si spread $>15\%$; $0.25$ ($>10\%$); $0.50$ ($>5\%$); $1.0$ ($\le 5\%$). \\
Régimen $G$ & $1.0$ si alineado; $0.50$ si neutral; $0.25$ si se opone (proxy: tendencia de SPY). \\
Volatilidad $V(r)$ & $1.0$ si falta IV rank; $0.25$ ($>80$); $0.50$ ($60$--$80$); $0.75$ ($<20$). \\
Concentración $K_{\text{conc}}$ & $1.0$ a $0.5$ según HHI. \\
Monte Carlo $K_{\text{MC}}$ & $0.5$--$1.0$ según CVaR/EV/PoP. \\
Sentimiento $K_{\text{sent}}$ & $0.75$ ante hype alto o desacuerdo confiado; $1.0$ en otro caso. \\
\bottomrule
\end{longtable}

Luego se aplica el \textbf{tope de Kelly} y los topes del portafolio:
\begin{align}
R_{\text{uncapped}} &= \min\big(R_{\text{mult}},\ R_{\text{Kelly}}\big), \\
\text{SuggestedRisk} &=
\begin{cases}
\min\!\big(R_{\text{uncapped}},\ \text{sleeve}_{\text{rest}},\ \text{mensual}_{\text{rest}},\ \text{abierto}_{\text{rest}}\big), & \text{si aprobada},\\
0, & \text{en otro caso}.
\end{cases}
\end{align}

\subsection{Límites por defecto (cuenta de 100{,}000 \MXN)}
\begin{longtable}{p{0.50\linewidth} p{0.30\linewidth}}
\toprule
\textbf{Parámetro} & \textbf{Valor por defecto} \\
\midrule
Valor total del portafolio & 100{,}000 \MXN \\
Sleeve de opciones ($3\%$) & 3{,}000 \MXN \\
Riesgo base por trade ($0.5\%$) & 500 \MXN \\
Pérdida mensual máxima ($1\%$) & 1000 \MXN \\
Riesgo abierto máximo ($2\%$) & 2{,}000 \MXN \\
\bottomrule
\end{longtable}

% ==========================================================
\section{Descubrimiento e ideas automáticas}
% ==========================================================

\subsection{Descubrimiento (\texttt{lib/discovery.ts})}
Combina varias fuentes (Yahoo \emph{day\_gainers/losers/most\_actives}, un
universo líquido de mega-caps y menciones de Reddit) y puntúa cada candidato con
componentes de configuración técnica, volatilidad realizada, fuerza relativa,
ratio de volumen y señal social, aplicando una \textbf{penalización por
extensión} cuando el movimiento intradía es grande (para no perseguir). Hay dos
perfiles de pesos: \emph{quality} (configuraciones limpias) y \emph{movers}
(reactivo a momentum y volumen).

\subsection{Ideas automáticas (\texttt{lib/auto-ideas.ts})}
\texttt{ensureDailyAutoIdeas} genera hasta cierto número de ideas por día:
descubre candidatos, arma el contexto de mercado, selecciona el mejor contrato de
opción vía Alpaca (\texttt{lib/option-selector.ts}), construye el
\texttt{TradeInput}, lo puntúa y lo guarda como \emph{watchlist} o \emph{reject}
(nunca aprobada automáticamente, conservando la aprobación manual).

\subsection{Selección de contrato (\texttt{lib/option-selector.ts})}
Puntúa cada contrato de la cadena con una rúbrica: ajuste de delta ($22\%$,
objetivo $0.40$), calidad de spread ($20\%$), liquidez ($18\%$, volumen + OI),
ajuste de DTE ($14\%$, máximo en $\sim 42$ días), IV ($10\%$), theta ($8\%$),
gamma ($4\%$) y \emph{moneyness} ($4\%$), con un umbral mínimo de $48/100$.

\subsection{Refresco diario (\texttt{lib/daily-refresh.ts})}
Orquesta la actualización: refresca cotizaciones de la watchlist, obtiene el
historial de SPY/QQQ/VIX, reconstruye el contexto de mercado, actualiza
técnicos y snapshots de Alpaca por idea, re-puntúa y guarda un
\texttt{TradeScoreSnapshot} para seguir la evolución diaria.

% ==========================================================
\section{Métricas de la bitácora}
% ==========================================================

En \texttt{lib/calculations.ts} se calculan, para operaciones cerradas:

\begin{align}
\text{WinRate} &= \frac{\#\text{ganadoras}}{\#\text{cerradas}}\times 100, \\
\text{Expectancy} &= \tfrac{\text{WinRate}}{100}\,\overline{\text{ganancia}} + \tfrac{\text{LossRate}}{100}\,\overline{\text{pérdida}}, \\
\text{Sharpe}_{\text{proxy}} &= \frac{\overline{P\&L}}{\sigma_{P\&L}}\sqrt{\frac{252}{\overline{\text{duración}}}}, \\
\text{Sortino}_{\text{proxy}} &= \frac{\overline{P\&L}}{\sigma_{\text{downside}}}\sqrt{\frac{252}{\overline{\text{duración}}}}.
\end{align}

Además calcula intervalos de confianza al $95\%$ por \emph{bootstrap} (generador
congruencial lineal con semilla determinista) y un \emph{backtest por bandas de
puntaje} ($<60$, $60$--$74$, $75$--$84$, $\ge 85$) para comparar el P\&L por
banda y validar si el modelo discrimina.

% ==========================================================
\section{Bot de investigación en Python}
% ==========================================================

\texttt{options-strategy-bot/} es un banco de pruebas independiente que
\textbf{no opera en vivo}. Replica la lógica de \RiskGate{} para escanear,
puntuar, hacer \emph{backtest} y llevar un libro de \emph{paper trading}.

\begin{longtable}{p{0.26\linewidth} p{0.66\linewidth}}
\toprule
\textbf{Módulo} & \textbf{Función} \\
\midrule
\texttt{config.py} & Esquema de configuración (riesgo, estrategia, costos, universo). \\
\texttt{data.py} & Carga de datos: CSV, Yahoo o barras sintéticas; modo \emph{offline}. \\
\texttt{features.py} & Indicadores (SMA, RSI, volatilidad realizada) y régimen de mercado. \\
\texttt{options.py} & Black--Scholes propio (precio y Greeks), proxy de IV y construcción de spreads de débito. \\
\texttt{model.py} & Puntaje de cuatro módulos, decisión y dimensionamiento estilo \RiskGate{}. \\
\texttt{strategies.py} & Genera candidatos por estrategia habilitada y los ordena por puntaje. \\
\texttt{backtest.py} & Backtest histórico con entradas/salidas, curva de equity y resumen (expectancy, win rate). \\
\texttt{paper.py} & Libro de paper trading en JSON. \\
\texttt{providers/alpaca.py} & Snapshot de cadena de opciones (solo lectura). \\
\texttt{cli.py} & Comandos \texttt{scan}, \texttt{backtest}, \texttt{paper-step}. \\
\bottomrule
\end{longtable}

\subsection{Black--Scholes y proxy de IV}
El bot valúa opciones con la forma cerrada de Black--Scholes (precio, delta,
gamma, theta, vega) y aproxima la IV a partir de la volatilidad realizada a 20
días ($\sigma_{\text{impl}}\approx 1.15\,\sigma_{\text{real}}$, acotada a
$[0.18, 0.85]$). Esto le permite simular el valor de estructuras sin acceso a
cadenas históricas reales.

\subsection{Scripts de validación}
\begin{itemize}[noitemsep]
  \item \texttt{backtest\_riskgate\_premiums.py}: lee contratos guardados en la
        base de \RiskGate{} y traza cómo se movió la prima real (Alpaca) tras la
        entrada.
  \item \texttt{compare\_scored\_ideas.py}: compara la precisión direccional y por
        bandas de puntaje de las ideas guardadas.
  \item \texttt{chart\_yahoo\_signal\_followthrough.py}: grafica el movimiento del
        subyacente tras una señal usando cierres ajustados de Yahoo.
\end{itemize}

% ==========================================================
\section{Limitaciones actuales}
% ==========================================================

\begin{itemize}[noitemsep]
  \item El régimen de mercado todavía puede depender de selección manual y usa la
        tendencia de SPY como \emph{proxy} de régimen para el multiplicador.
  \item La probabilidad de ganancia y el tope de Kelly dependen de IV y datos de
        contrato; cuando faltan, se usan \emph{proxies} (probabilidad implícita
        por puntaje).
  \item Yahoo Finance es una fuente pública no oficial; puede fallar, retrasarse o
        cambiar de formato.
  \item No hay detección automática de fechas de reportes (earnings) ni de eventos
        macro programados.
  \item El sentimiento de Reddit usa TF--IDF local, no un modelo de lenguaje, por
        lo que capta coincidencia de términos más que matices.
  \item El backtest no modela el deslizamiento real de ejecución ni el llenado
        parcial de spreads ilíquidos con fidelidad de mercado.
  \item No hay autenticación multiusuario ni cifrado de la base local.
\end{itemize}

% ==========================================================
\section{Recomendaciones y próximos pasos}
% ==========================================================

Esta sección reúne las mejoras propuestas, organizadas por horizonte de
implementación. Se incluye explícitamente la mejora de los \emph{proxies} que el
modelo usa hoy.

\subsection{Corto plazo: robustez de datos y \emph{proxies}}

\paragraph{Reemplazar el proxy de régimen de SPY.}
Hoy el multiplicador de régimen $G$ usa únicamente la tendencia de SPY como
proxy. Se recomienda construir un \emph{régimen compuesto} que combine SPY, QQQ,
breadth y VIX en un solo índice $\in[-1,1]$ y derivar $G$ de él, de modo que un
mercado con SPY plano pero breadth débil y VIX en alza no reciba $G=0.5$.

\paragraph{Sustituir la probabilidad implícita por puntaje en Kelly.}
Cuando faltan IV o datos del contrato, el tope de Kelly usa una probabilidad
derivada del puntaje ($p=S/100$). Conviene migrar a una PoP calculada siempre
desde la cadena real (delta como proxy de probabilidad o PoP lognormal completa),
y registrar en la idea si la PoP fue \emph{medida} o \emph{estimada}.

\paragraph{Mejorar el proxy de IV del bot.}
El factor fijo $1.15\times\sigma_{\text{real}}$ podría calibrarse por subyacente
y por régimen (term structure y \emph{skew} simplificados), o sustituirse por IV
real cuando exista snapshot de Alpaca.

\paragraph{Endurecer la fuente de datos.}
Añadir reintentos, validación de esquema y \emph{fallback} entre proveedores
(Yahoo $\rightarrow$ Alpaca $\rightarrow$ manual) con marcas de frescura visibles
en la interfaz, y registrar cuándo un dato es estimado.

\subsection{Mediano plazo: cobertura del modelo}

\begin{itemize}[noitemsep]
  \item \textbf{Calendario de earnings y eventos macro}: marcar ideas cuyo
        vencimiento cruce un reporte o un evento (FOMC, CPI) y aplicar un tope o
        advertencia automática.
  \item \textbf{Checks por estrategia}: validaciones específicas para
        \emph{long calls/puts}, \emph{debit spreads}, \emph{protective puts} y
        \emph{collars} (ancho, break-even, reward/risk, cobertura).
  \item \textbf{Correlación entre ideas}: extender la concentración por sector a
        una matriz de correlación entre subyacentes para evitar riesgo redundante.
  \item \textbf{Calibración del modelo con la bitácora}: usar el backtest por
        bandas de puntaje para re-ajustar pesos y umbrales (validación walk-forward).
  \item \textbf{Sentimiento con embeddings semánticos}: reemplazar TF--IDF por
        \emph{embeddings} de un modelo de lenguaje local para captar matices y
        sarcasmo, manteniendo el procesamiento en la máquina.
\end{itemize}

\subsection{Largo plazo: producto y plataforma}

\begin{itemize}[noitemsep]
  \item \textbf{Empaquetado de escritorio}: envolver la app con Tauri para generar
        un ejecutable de Windows ligero, manteniendo la primera versión de
        escritorio en modo solo lectura y sin credenciales del broker.
  \item \textbf{Cifrado de la base local} y bloqueo por contraseña local.
  \item \textbf{Auditoría}: bitácora inmutable de decisiones del modelo para
        revisar por qué se aprobó o rechazó cada idea.
  \item \textbf{Integración de datos profesional}: migrar de Yahoo a un proveedor
        formal de cadenas y Greeks.
\end{itemize}

\subsection{Candidatos de integración externa}

\begin{longtable}{p{0.20\linewidth} p{0.58\linewidth}}
\toprule
\textbf{Proveedor} & \textbf{Uso recomendado} \\
\midrule
Tradier & Cadenas de opciones con Greeks/IV, datos de cuenta y sandbox; ruta hacia paper/live. \\
Polygon.io & Datos de mercado de opciones de EE.\,UU.\ (tiempo real e histórico). \\
MarketData.app & Vencimientos, strikes, cadenas y cotizaciones para prellenado. \\
Alpaca & Cadena de opciones ya integrada; paper trading futuro con confirmación explícita. \\
\bottomrule
\end{longtable}

\begin{itemize}[noitemsep]
  \item Tradier: \url{https://docs.tradier.com/reference/brokerage-api-markets-get-options-chains}
  \item Polygon.io: \url{https://polygon.io/docs/rest/options/overview}
  \item MarketData.app: \url{https://www.marketdata.app/docs/api/}
  \item Alpaca: \url{https://docs.alpaca.markets/docs/options-trading}
\end{itemize}

\noindent
La ruta recomendada es siempre \textbf{datos de solo lectura primero}. Cualquier
automatización de órdenes debe quedar fuera de alcance hasta que existan manejo
confiable de datos, pruebas, auditoría y flujos de confirmación explícita.

% ==========================================================
\section{Cómo ejecutar}
% ==========================================================

\subsection{Aplicación web}
\begin{lstlisting}
npm install
npm run prisma:generate
npm run db:init
npm run seed
npm run dev
\end{lstlisting}
Abrir \texttt{http://127.0.0.1:3000}. Para datos de opciones, configurar
\texttt{APCA\_API\_KEY\_ID} y \texttt{APCA\_API\_SECRET\_KEY} en \texttt{.env}.

\subsection{Bot de investigación}
\begin{lstlisting}
cd options-strategy-bot
pip install -r requirements.txt
python -m bot.cli scan --synthetic --offline
python -m bot.cli backtest --synthetic
\end{lstlisting}

% ==========================================================
\section{Notas de seguridad}
% ==========================================================

\begin{itemize}[noitemsep]
  \item Las opciones son riesgosas y pueden expirar sin valor.
  \item \RiskGate{} es una herramienta de apoyo a la decisión, no asesoría
        financiera.
  \item La aplicación no debe almacenar contraseñas del broker.
  \item Las integraciones de solo lectura deben preceder a cualquier flujo de
        colocación de órdenes.
  \item El modelo es cuantitativo y basado en reglas, no un sistema predictivo
        garantizado.
\end{itemize}

% ==========================================================
\section{Conclusión}
% ==========================================================

\RiskGate{} reúne los componentes de una disciplina de opciones rigurosa:
aislamiento por portafolio, posiciones manuales e importadas, watchlist,
descubrimiento de candidatos, una compuerta formal de ideas con puntaje
ponderado adaptativo, calidad de opciones continua, analítica probabilística
(PoP, EV, Kelly, VaR/CVaR), simulación Monte Carlo, sentimiento opcional de
Reddit, control de concentración, dimensionamiento acotado por la cuenta y una
bitácora con métricas de validación. El bot de Python complementa la app como
banco de pruebas de \emph{backtesting}. La evolución más valiosa es continuar
\textbf{automatizando inputs objetivos y mejorando los proxies del modelo}
mientras se conserva la aprobación manual y el principio local-first.

\end{document}
